% ============================================================================
% BOLUM 5: TARTISMA
% ============================================================================
\section{Tartışma}
\label{sec:discussion}

% ----------------------------------------------------------------------------
% 5.1 Ana Bulgular
% ----------------------------------------------------------------------------
\subsection{Ölçümlerin Gösterdikleri}
\label{subsec:findings}

\paragraph{Protokol varyansı, karşılaştırılan büyüklüklere hükmediyor}

Merkezî sonuç bir yayılım sıralamasıdır. Transfer asimetrisinin kendisi üzerinde ölçüldüğünde, her iki mimariyi sıfırdan yeniden eğitmek onu $0{,}23$--$1{,}48$ puanlık standart sapmayla oynatırken, yalnızca koşullama protokolünü değiştirmek $10{,}45\pm0{,}76$ puan oynatmaktadır (Bölüm~\ref{subsec:statistics_results}). İki yayılımı kat değeri olarak değil mutlak birimlerle veriyoruz; çünkü üç koşumla o oranın aralığı bir manşeti taşıyamayacak kadar geniştir. Farklı bir nicelik üzerinde ölçek fikri vermek üzere: bir saldırıyı yeni bir rastgele başlangıçla yeniden koşmak mutlak gürbüz doğruluğu yaklaşık $0{,}05$ puan, yeniden eğitmek ise çekişmeli eğitilmiş çifti en çok $0{,}55$ puan oynatır. Ölçüm seçiminin on puan oynattığı, yeniden eğitmenin ise bir puan oynattığı bir büyüklük mimarilerin kararlı bir özelliği değildir ve protokol olmadan raporlandığında neredeyse hiçbir şey iletmez. Bu, koşullamaya karşı bir argüman değildir --- doğru sınıflandırılan örneklerle koşullamak yerleşik pratiktir~\cite{liu2017delving, dong2018boosting, ravikumar2023trend} ve gerekli olmaya devam etmektedir --- seçimin dipnota değil sonuca ait olduğuna dair bir argümandır.

Dört protokol birbirinin yerine geçmez ve her birinin savunulabilir bir kullanımı vardır. Koşulsuz oran mimari karşılaştırması için hiç kullanılmamalıdır: $3\times3$ matrisimiz, koşullu orandan sapmasının hedefin temiz hatasının $(1 - r_{\text{koş}})$ ile ölçeklenmiş hâline cebirsel olarak eşit olduğunu gösterir; yani transfer edilebilirliği değil, bilinen bir oranda hedef zayıflığını ölçer. Hedef doğru protokolü bu karıştırıcıyı giderir ama kaynağın kendisinin çözemediği örnekleri kabul eder; bu marjinal altkümeler aşırı kırılgan ve asimetrik boyutlu olduğundan iki yönü, temiz doğruluk farkına bağlı bir miktarda birbirine doğru çeker. Her ikisi doğru protokolü iki yönü özdeş örnekler üzerinde puanlar ve dördü içinde eşleştirilmiş çıkarımı destekleyen tek protokoldür; birincil tahmin olarak onu bu yüzden alıyoruz. Başarılı kaynak protokolü farklı bir soruyu yanıtlar --- zaten işe yaramış bir saldırı ne kadar taşınır --- ve daha büyük değeri, kaynak başarısıyla koşullamanın güçlü pertürbasyonları seçmesini yansıtır.

\paragraph{Yön yalnızca modeller eşitlenmemişken kararlıdır}

Dört protokol, üç eğitim koşusu ve iki saldırı ailesi (PGD-10 ve MI-FGSM) boyunca, CNN'de üretilen çekişmeli örnekler \emph{her iki CIFAR veri kümesinde} ViT'e tersinden daha iyi transfer olmaktadır. Aynı tarifle çekişmeli eğitimin iki mimariyi $11$ ya da $21$ puan değil yalnızca $1{,}85$ puan ayırdığı SVHN'de ise bu artık doğru değildir: iki protokol CNN üstünlüğü, iki protokol ViT üstünlüğü bildirmekte ve birincil her ikisi doğru protokolünde fark $-0{,}35$ puan olup eşdeğerlik kabul edilmektedir (Bölüm~\ref{subsec:svhn}). Yön, mimarilerin değil eşitlenmemiş rejimin bir özelliğidir. Değişen, etkinin ne kadar büyük göründüğü ve bir eşdeğerlik testinin onu kabul edip etmediğidir. Üç hedefli analizimiz bu yön için CNN pertürbasyonlarının özel olmasını gerektirmeyen bir okuma önermektedir: gelen transfer her hedefin kendi beyaz kutu kırılganlığını izler (üç hedef üzerinde $r = 0{,}986$) ve en gürbüz hedef iki aileden gelen saldırıları neredeyse eşit soğurur. Bu okumaya göre asimetri, hangi modelin daha zayıf hedef olduğuna dair bir ifadedir; dönüştürücüye özgü bir saldırının~\cite{zhang2023tgr} ViT yönünü kurtaramaması da bununla tutarlıdır. Üç hedefle korelasyon yönlendiricidir; kesinleştirmek daha büyük bir model havuzu gerektirir.

\paragraph{Manşet sayının içindeki atıf, manşetin kendisinden daha kararsız}

Koşullu ayrıştırma, mutlak gürbüzlüğü temiz doğruluk çarpanı ile koşullu duyarlılık çarpanına ayırır. Düzeltilmiş koşularımızda ikisi de CNN lehinedir; ayrıştırma ile manşet uyuşur. Önceki kontrol noktalarımızda uyuşmuyorlardı: koşullu duyarlılık PGD-10 altında eşleşik, AutoAttack altında hafifçe ViT lehine görünüyordu. Mutlak sıralama iki kontrol noktası kümesinde de özdeş olduğundan, yalnızca mutlak doğruluğu raporlayan bir makale hiçbir değişiklik görmezdi; tek koşudan ayrıştırma raporlayan bir makale ise tekrarlanmayan mekanizma düzeyinde bir sonuç çıkarırdı. Pratik ders şudur: mekanizma iddiaları sıralama iddialarından daha çok eğitim koşusu ister ve ikisine aynı kanıt ağırlığı verilmemelidir.

\paragraph{Ölçümden sağ çıkamayan üç iddia}

Seyreklik mekânsal lokalite değildir. Çekişmeli eğitilmiş CNN gradyanları kütlenin daha büyük kısmını daha az bileşene yoğunlaştırır (Hoyer $0{,}4928$'e karşı $0{,}4561$, eşleştirilmiş $p<10^{-11}$); ama aynı gradyanların dört mekânsal betimleyicisi, bu bileşenlerin görüntüdeki dizilişinde karşılık gelen bir fark göstermez. CNN gradyanlarının mekânsal olarak lokalize olduğu betimlemeleri bu yüzden ayrı kanıt ister; seyreklik ölçüleri~\cite{hurley2009comparing} bunu sağlamaz.

Çekişmeli pertürbasyon CLS dikkat entropisini ölçülebilir biçimde değiştirmez. $n=1000$'de katman başına entropi değişimi en çok $0{,}005$ nattır; dikkat haritaları ise küçük ve derinlikle artan bir miktarda yer değiştirir. Yama saldırılarıyla elde edilen dikkat kayması sonuçları~\cite{gu2022vision, fu2022patchfool} farklı bir tehdit modeline aittir ve ölçüm olmadan $L_\infty$ sınırlı pertürbasyonlara genellenmemelidir.

ViT'e özgü bir transfer saldırısı burada ViT kaynaklı transferi iyileştirmez. TGR~\cite{zhang2023tgr} 14,9 puan beyaz kutu gücüne mal olmakta ve bunun hiçbirini hedefte geri vermemektedir. Bunu rejime özgü bir negatif sonuç olarak raporluyoruz --- hedeflerimiz çekişmeli eğitilmiş, verimiz CIFAR ölçeğinde --- ve tam da bu türden tek bir test edilmemiş varsayımın, protokole duyarlı bir karşılaştırmaya taşıyacağı yük yüzünden.

\paragraph{Ayakta kalanlar}

İki davranışsal fark bağımsız kontrol noktası kümelerinde tekrarlanmakta ve protokol sorusundan etkilenmemektedir. Gradyan yapısı farklıdır: CNN girdi gradyanları üç ölçekten bağımsız ölçünün üçünde de daha seyrektir, bu sıralamayı mimari değil çekişmeli eğitim yaratır (temiz eğitilmiş çiftte tersidir) ve örnekler arası hizalanma ViT'te mutlak kosinüste daha yüksekken işaretli ortalama iki modelde de sıfıra yakındır --- örnekler yönleri değil duyarlılık eksenlerini paylaşır; ölçümümüzü bir evrensel pertürbasyon iddiasından~\cite{moosavi2017universal} ayıran ayrım budur. Derinlik profilleri farklıdır: CNN'de çekişmeli hasar son artık aşamaya sıkışmış geç, yerel, büyüklük değiştiren bir olay; ViT'te dağıtık, yalnızca yön değiştiren bir dönmedir. ViT'in ölçülen profilinin jeton agregasyonuna bağlı olması --- CLS yolunun yama ortalamasının üç katı kayması --- katman bazlı çalışmalar için kendi başına bir raporlama gerekliliğidir.

% ----------------------------------------------------------------------------
% 5.2 Oneriler
% ----------------------------------------------------------------------------
\subsection{Raporlama Gereklilikleri}
\label{subsec:implications}

Yukarıdaki ölçümler, mimariler arası gürbüzlük çalışmaları için kısa ve somut bir liste doğurmaktadır.

\begin{enumerate}
\item \textbf{Koşullama protokolünü belirtin ve eşleştirilmiş olanları yeğleyin.} Paydaya hangi örneklerin girdiğini raporlayın. Yönlü bir iddia yapılacaksa, eşleştirilmiş çıkarım mümkün olsun diye iki yönü özdeş örnekler üzerinde puanlayan bir protokol kullanın.
\item \textbf{Koşulsuz transfer oranlarını mimari karşılaştırması olarak raporlamayın.} Hedefler temiz doğrulukta ayrışıyorsa ham oran, transfer edilebilirlik ile hedef zayıflığının \emph{bilinen} bir karışımıdır: $r_{\text{ham}} - r_{\text{koş}} = e\,(1 - r_{\text{koş}})$. Bu uydurulmuş bir ilişki değil bir özdeşlik olduğundan, düzeltme yayımlanmış ham oranlara hiçbir şey yeniden koşulmadan uygulanabilir.
\item \textbf{Yalnızca skoru değil, ayrıştırmayı raporlayın.} Mutlak gürbüzlük, temiz doğruluk ile koşullu duyarlılığın çarpımına birebir ayrışır; manşeti hangi çarpanın taşıdığı işin ilginç kısmıdır ve skordan geri çıkarılamaz.
\item \textbf{Sıralama iddialarını mekanizma iddialarından ayırın.} Sıralamalarımız kontrol noktası kümeleri arasında kararlıydı; mekanizma atfımız değildi. Mekanizma iddiaları daha çok eğitim koşusu taşımalıdır.
\item \textbf{Katman bazlı analizlerde agregasyonu belirtin.} Aynı blokların CLS, yama ortalaması ve tüm jeton özetleri farklı kayma büyüklükleri ve farklı minimumlar verir.
\item \textbf{Dikkat ve lokalite iddialarını gerçek bir örneklem boyutunda sınayın.} Raporladığımız iki sıfır sonuç da küçük bir figür partisinde pozitif görünürdü.
\end{enumerate}

Bu kıyaslamada eşleşmiş standart çekişmeli eğitim altında iki aile arasında seçim yapan uygulayıcılar için: CNN daha yüksek mutlak gürbüz doğruluk vermekte ve üstünlük ölçtüğümüz her görünümde korunmaktadır. \%62{,}76'daki WideResNet referansı, tarif, kapasite ve ek verinin birlikte~\cite{gowal2020uncovering} gürbüzlüğü mimari seçiminden daha ileri taşıdığını göstermektedir; gürbüzlük stratejisi olarak mimari seçimine yatırım yapmadan önce tartılmaya değer.

% ----------------------------------------------------------------------------
% 5.3 Sinirliliklar
% ----------------------------------------------------------------------------
\subsection{Sınırlılıklar}
\label{subsec:limitations}

Çalışmamız bilinçli olarak dardır ve bu darlık sonuçları belirli biçimlerde sınırlar.

Protokol duyarlılığı sonucu üç veri kümesi ve tek model çifti üzerinde ölçülmüştür. CIFAR-100'de yayılım CIFAR-10'dakinden büyüktür ($13{,}58\pm1{,}71$'e karşı $10{,}45\pm0{,}76$ puan) ve işaret on iki ölçümün tamamında korunmaktadır (Bölüm~\ref{subsec:cifar100}); bu, CIFAR'a özgü bir şeyden değil hedefler arasındaki temiz doğruluk farkından geçen bir mekanizmayla tutarlıdır. Ancak bu açıklama eksiktir ve sloganı değil sınırı yazıyoruz: dört protokolün üçü için geçerlidir; başarılı kaynak protokolü ise \emph{kaynağın} gürbüzlüğüne bağlı ikinci bir sürücü ekler ve bu sürücü genel ilişkiyi tersine çevirecek kadar güçlüdür (Bölüm~\ref{subsec:transfer}). Her iki bağımlılığın işlevsel biçimini de güvenle belirtemiyoruz: kestirimler altı tohum-ve-veri-kümesi kümesine dayanmaktadır ve havuzumuzdaki küçük farklı çiftlerin tamamı harici WideResNet'i içerdiğinden, küçük temiz doğruluk farkı ile farklı eğitim tarifi burada birbirinden ayrılamamaktadır. Değişen temiz doğruluk farklarına sahip birden çok mimari çiftini tarayan bir çalışma, bu ölçümleri bir kalibrasyon eğrisine dönüştürecektir.

İki kontrol noktası kümesi arasındaki farkların nedenini yalıtamıyoruz. Düzeltilmiş protokol, doğrulama düzenini değiştirmiş ve aynı anda yeni bir eğitim koşusu üretmiştir. Veri miktarı açıklaması dışlanmıştır (düzeltilmiş temiz modeller daha az veriyle daha iyidir); ancak seçim sızıntısını olağan koşudan koşuya değişimden ayırmak, özdeş veriyle sızıntılı ve sızıntısız seçim altında eğiten bir ablasyon gerektirir. Karşıtlığı sızıntının ölçümü olarak değil, tek koşuluk mekanizma iddialarının kırılganlığına dair kanıt olarak raporluyoruz.

Bu makaledeki her ölçüm tek bir tehdit modelinin içinde yapılmıştır: $\eps = 8/255$ ile $\Linf$. Protokol yayılımının norm değişince ayakta kalıp kalmadığı ayrı bir sorudur ve kaydedilmiş ama henüz yanıtlanmamıştır. Böyle bir sınamanın neyi gösterip neyi gösteremeyeceğini belirtmek gerekir: modellerimiz $\Linf$ ile eğitilmiştir, dolayısıyla onları bir $L_2$ bütçesi altında değerlendirmek $L_2$ gürbüzlüğünü değil ölçüm eksenini yoklar ve çıkan mutlak değerler $L_2$ ile eğitilmiş referanslarla karşılaştırılamaz.

Mimari başına üç koşu tohum varyansını sınırlar ama tarif uzayını taramaz; küçük veri kümelerinde ViT çekişmeli eğitiminin tarife duyarlı olduğu bilinmektedir~\cite{mo2022adversarial, debenedetti2023light} ve difüzyonla üretilmiş sentetik veriyle ViT'lerin CIFAR-10'da WideResNet'leri geçtiği bildirilmiştir~\cite{wang2023better, wu2025vision}. ViT sonuçlarımız ulaşılabilir en iyi ViT gürbüzlüğünü değil, temel çekişmeli eğitimi betimler. ViT-Tiny hattı $32\times32$ girdileri model içinde $224\times224$'e büyütür; CIFAR'a özgü kontrolümüz gradyan sıralamasının büyütme artefaktı olmadığını gösterir, ama bu kontrol kapasite ve yama boyutunda farklıdır; dolayısıyla karıştırıcıyı kalibre eder, ikinci bir eşleşmiş çift sağlamaz. İki model parametre sayısında da ayrışır (11,2M'ye karşı 5,7M); mutlak farkın mimariye atfı bu yüzden kısmi kalır. Son olarak tek bütçeli bir $L_\infty$ tehdit modeli değerlendiriyoruz ve TGR sonucumuz, yöntemin yayımlanmış kurulumunun yeniden üretimi değil bir uyarlamasıdır.

% ----------------------------------------------------------------------------
% 5.4 Gelecek Calismalar
% ----------------------------------------------------------------------------
\subsection{Gelecek Yönelimler}
\label{subsec:future}

En doğrudan uzantı, protokol duyarlılığı ölçümünü genel bir ölçüme dönüştürmektir: kontrollü temiz doğruluk farklarına sahip mimari çiftlerini tarayıp protokoller arası yayılımın bu farkla nasıl ölçeklendiğini ölçmek. İlişki, üç hedefli korelasyonumuzun önerdiği kadar düzenliyse, yayımlanmış ham oran asimetrilerinin atılmak yerine yeniden yorumlanmasına imkân veren bir düzeltme sağlar.

Daha önce gelecek çalışma diye sıraladığımız uzantılardan ikisi bu arada koşulmuştur; hem neyi karara bağladıklarını hem neyi bağlamadıklarını kaydediyoruz. Sızıntı ablasyonu altı tam bütçeli yörünge üzerinde yürütülmüş ve sızıntının kontrol noktası seçimine etkisini \emph{saptamamıştır}; başlangıçta konan sorunun daha zayıf bir sürümünü yanıtlar, çünkü tedavi tek bir doğrulama bölmesine uygulanmıştır ve o tasarımda bölme kimliği ile sızıntı birbirinden ayrıştırılamaz. Bilgi verici yan ürünü, Bölüm~\ref{subsec:statistics_results}'te raporlanan seçim protokolü yayılımıdır. İkinci veri kümesi de koşulmuştur (Bölüm~\ref{subsec:cifar100}). Açık kalan, kapasitesi eşlenmiş çifttir (örneğin ViT-Small'a karşı ResNet-50); bu, ayakta kalan davranışsal farkların --- seyreklik sıralaması ve derinlik profilleri --- ailelerin mi yoksa bu küçük modellerin mi özelliği olduğunu test edecektir.
